\input{../universal-pre}

\title{\tbf{DGH2021}}
\author{\tbf{Notes} by Cong Wen, 03/2022}
\date{}

\begin{document}
\maketitle
\thispagestyle{empty}
\begin{abstract}
  This is my self-use notes\footnote{Last updated on March 10, 2022} on \tbf{Uniform Mordell-Lang}\cite{DGH2021}. Please reach out if you find anything incorrect.
\end{abstract}
\tableofcontents


% ----------------------------------------------------------------------------------------------------


\section{introduction}
\begin{stp}
  Let $F$ be any field, a curve on $F$ is a geometrically irreducible, projective variety of dimension $1$ defined over $F$.

  $C$ is a smooth curve of genus $g\geq 2$ defined over a number field $F$, then $\sP{C(F)}<+\infty$\cite{Faltings1983}. $\Jac(C)$ the Jacobian variety of $C$, then $\rho=\Rank\Jac(C)(F)<+\infty$ (Mordell-Weil).
\end{stp}

\begin{thm}
  Let $F$ be a number field with $[F:Q]=d\geq 1$, $C$ smooth curve on $F$ of genus $g\geq 2$, then there exists $c(g,d)$,
  \[
    \sP{C(F)}\leq c(g,d)^{1+\rho}.
  \]
\end{thm}
\begin{rmk}
  This is the main theoreom of the paper. The main innovation is the proof of a new bound for N{\'e}ron-Tate height.

  The following two theorems are further results with this approach.
\end{rmk}


\begin{stp}
  Let $\bA_{g,1}$ be the coarse moduli space of principally polarized abelian varieties of dimension $g$ with $\iota:\bA_{g,1}\ar\bP_{\oL{\bQ}}^{m}$, fix a Weil height $h$ on $\bP_{\oL{\bQ}}^{m}$, denote the class $[\Jac(C)]\in\bA_{g,1}(\oL{\bQ})$.
\end{stp}

\begin{thm}
  Take any subgroup $\Gamma\subset\Jac(C)(\oL{\bQ})$ of finite rank $\rho\geq 0$, there exists $c_{1}(g,\iota)\geq 0, c_{2}(g,\iota)\geq 1$ such that if $h\sR{\iota([\Jac(C)])}\geq c_{1}$, then
  \[
    \sP{C(\oL{\bQ})-P_{0}}\cap\Gamma\leq c_{2}^{1+\rho}.
  \]
\end{thm}

\begin{cor}
  Applying to $\Gamma=\Jac(C)(\oL{\bQ})_{\mathrm{tor}}$, we have
  \[
    \sP{C(\oL{\bQ})-P_{0}}\cap\Jac(C)(\oL{\bQ})_{\mathrm{tor}}\leq c_{2}.
  \]
\end{cor}

\begin{thm}
  Let $F$ be a number field with $[F:\bQ]=d\geq 1$, $C$ smooth curve on $F$ of genus $g\geq 2$, then there exists $c(g,d)$, for any $P_{0}\in C(\oL{\bQ})$,
  \[
    \sP{C(\oL{\bQ})-P_{0}}\cap\Jac(C)(\oL{\bQ})_{\mathrm{tor}}\leq c(g,d).
  \]
\end{thm}
\begin{rmk}
  \red{Why working over number fields drop the condition on height?}

  The following is the strategy of the paper.
\end{rmk}

\begin{stp}
  Let $k=\oL{\bQ}$, $S$ regular irreducible, quasi-projective, variety defined over $k$ which is Zariski open in irreducible projective variety $\oL{S}\subset\bP_{k}^{n}$. Let $\pi:\cA\ar S$ be an abelian scheme of relative dimension $g\geq 1$ with a closed immersion $\cA\ar\bP_{k}^{n}\times S$ over $S$, also assume the closed immersion on the generic fiber $A=\cA_{\eta}\ar\bP_{k{S}}^{n}$ is given by $l\cL, l\geq 4$, $\cL$ symmetric ample.

  Now we have the usual height $h:\oL{S}(k)\ar\bR$ and N{\'e}ron-Tate(canonical) height $\oH{h_{\cA}}:\cA(\oL{\bQ})\ar\bR$.

  Now let $\cA\ar S$ be an abelian scheme carrying a principal polarization and level $\ell\geq 3$ structure.
\end{stp}
\begin{dfn}
  An irreducible subvariety $X\subset\cA$ is said to be nondegenerate if there is an open non-empty $\Delta\subset S^{\mathrm{an}}$, with Betti map $b_{\Delta}:\cA_{\Delta}=\pi\xI(\Delta)\ar\bT^{2g}$, such that
  \[
    \max_{x\in X^{\mathrm{sm,an}}\cap\cA_{\Delta}}\Rank_{\bR}(\dif b_{\Delta}|_{X^{\mathrm{sm,an}}})_{x}=2\Dim X.
  \]
\end{dfn}

\begin{thm}\label{thm:htieqlt}
  Let $X\subset \cA\ar S$ be a closed irreducible subvariety dominating $S$, and $X$ nondegenerate. Then there exists $c_{1}>0, c_{2}\geq 0$, and a Zariski open $U\subset X$ such that
  \[
    \oH{h_{\cA}}(P)\geq c_{1}h(\pi(P))-c_{2}, \forall P\in U(\oL{\bQ}).
  \]
\end{thm}

\begin{fct}
  The moduli space $\bA_{g,\ell}$ of principally polarized abelian varieties of dimension $g$ and level-$\ell$-structure($\ell\geq 4$) is an irreducible regular quasi-projective variety defined over a number field, the moduli space $\bM_{g,\ell}$ of smooth curves of genus $g$ whose Jacobian is equipped with level-$\ell$-structure is also an irreducible regular quasi-projective variety defined over a number field.
\end{fct}


\section{Betti map and Bett form}

The existence and uniqueness (up to $\GL_{2g}(\bZ)$) of Bette maps and Betti forms.




\section{setup and notation for the height inequality}
\begin{stp}
  $k$ algebraically closed subfield of $\bC$,

  $S$ be a regular irreducible quasi-projective variety over $k$ which is Zariski open in an irreducible projective variety $\oL{S}\subset\bP_{k}^{m}$.

  $\pi:\cA\ar S$ abelian scheme with closed immersion $\cA\ar\bP_{k}^{n}\times S$ over $S$. Hence there is a closed immersion of generic fiber $A\ar S$ to $\bP_{k(S)}^{n}$, say its given by global sections of $lL$ with $L$ symmetric ample line bundle with $l\geq 4$. Also assume $\cA\ar S$ admits a symplectic level-$\ell$-structure.

  In particular, $S$ is the moduli space of principally polarized abelian varieties and $\cA$ is the universally family.
\end{stp}




\section{intersection theory and height inequality on the total space}

\begin{stp}
  Continue with the closed immersion $\cA\ar\bP_{k}^{n}\times S$, assume $k=\oL{\bQ}$. Let $X$ be a closed irreducible subvariety of $\cA$ of dimension $d$ defined over $\oL{\bQ}$ such that $\pi|_{X}:X\ar S$ is dominant, $\omega$ the Betti form on $\cA$.
\end{stp}

This whole section is devoted to the following height inequality
\begin{prp}\label{prp:htieqlt}
  Keep the setup above, suppose $X^{\mathrm{an}}$ contains a smooth point at which $\omega|_{X^{\mathrm{sm,an}}}^{\wedge d}>0$. Then there exists $c_{1}>0$ such that any $N\in\bN$ a power of $2$, $X$ has a Zariski dense open subset $U_{N}$ defined over $oL{\bQ}$ and a constant $c_{2}(N)$ such that
  \[
    h([N]P)\geq c_{1}N^{2}h(\pi(P))-c_{2}(N), \forall P\in U_{N}(\oL{\bQ}).
  \]
\end{prp}



\section{proof of the height inequality}

The proof directly comes from Theorem~\ref{thm:htieqlt} using Proposition~\ref{prp:htieqlt} and Silverman-Tate theorem.

\begin{thm}[Silverman-Tate]
  There exists a constant $c_{0}\geq 0$ such that
  \[
    \sP{\oH{h_{\cA}}(P)-h(P)}\leq c_{0}\max\{1, h(\pi(P))\}\leq c_{0}(1+h(\pi(P))), \forall P\in\cA(\oL{\bQ}).
  \]
\end{thm}


\section{prepartion for counting points}

This section fixes the setup for next chapter.
\begin{stp}
  Fix $g\geq 2$ integer, $\bM_{g}$ the fine moduli space of smooth curves of genus $g$ with level-$\ell$-structure and $\ell\geq 3$. Now $\bM_{g}$ is a regular, quasi-projective variety of dimension $3g-3$. The universal curve $\fC_{g}$ is smooth proper over $\bM_{g}$ with smooth fibers.
\end{stp}


Scholze~\autoref{prp:htieqlt}



\section{N{\'e}ron-Tate distance between points on curves}






\section{proof of Theorem~1.1, 1.2, and 1,4}






% ----------------------------------------------------------------------------------------------------
\bibliographystyle{alpha}
\bibliography{../universal-ref}
\end{document}
